%-----------------------------------------------------------------------------%
\chapter{\babEmpat}

Bab ini membahas eksperimen untuk menguji performa model klasifikasi EEG dengan augmentasi data berbasis \textit{Variational Autoencoder} (VAE). Eksperimen difokuskan pada analisis kualitas data sintetis yang dihasilkan serta pengaruhnya terhadap performa klasifikasi Autism Spectrum Disorder (ASD). Pembahasan mencakup spesifikasi sistem yang digunakan, skenario pengujian, proses augmentasi data, serta metrik evaluasi untuk menilai kualitas data sintetis dan performa model klasifikasi.

%-----------------------------------------------------------------------------%
\section{Lingkungan Training}
Bagian ini menjelaskan spesifikasi perangkat keras dan perangkat lunak yang digunakan dalam proses preprocessing sinyal EEG, training model VAE, generasi data sintetis, serta pelatihan dan evaluasi model klasifikasi.

\subsection{Spesifikasi Perangkat Keras}
Eksperimen dilakukan pada perangkat pribadi yang dilengkapi dengan RAM 16 GB berjenis DDR4 3200 MT/s dengan konfigurasi 2 $\times$ 8 GB. Proses pelatihan model diakselerasi menggunakan GPU NVIDIA GeForce RTX 3050 Laptop GPU dengan VRAM 4 GB, yang memungkinkan eksekusi paralel pada operasi tensor selama proses training model deep learning dan pemrosesan sinyal EEG. Sistem menggunakan prosesor Intel Core i5-12500H generasi ke-12 dengan 12 core dan 16 thread pada frekuensi dasar 3.10 GHz. Spesifikasi lengkap perangkat keras eksperimen disajikan pada Tabel \ref{tab:hardware-spec}.\begin{table}[H]
\centering
\footnotesize
\caption{Spesifikasi Perangkat Keras Eksperimen}
\label{tab:hardware-spec}
\begin{tabular}{|c|c|}
\hline
\textbf{Komponen} & \textbf{Spesifikasi} \\
\hline
CPU & Intel Core i5-12500H (12 Cores, 16 Threads, @ 3.10 GHz) \\
\hline
GPU & NVIDIA GeForce RTX 3050 Laptop GPU (4 GB VRAM) \\
\hline
RAM & 16 GB (2 $\times$ 8 GB) DDR4 3200 MT/s \\
\hline
Arsitektur & x86\_64 \\
\hline
\end{tabular}
\end{table} 
\subsection{Spesifikasi Perangkat Lunak}
Untuk lingkungan perangkat lunak, sistem dibangun di atas Python 3.9 dengan memanfaatkan TensorFlow (v2.10.0) sebagai framework utama dalam implementasi model \textit{Variational Autoencoder} (VAE). Selain itu, Scikit-learn (v1.6.1), Feature-engine (v1.9.4), dan XGBoost (v2.1.4) digunakan dalam proses preprocessing, seleksi fitur, serta klasifikasi data EEG. Proses pengolahan sinyal dan ekstraksi fitur didukung oleh SciPy (v1.13.1), sedangkan pengolahan data numerik dan visualisasi dilakukan menggunakan NumPy (v1.23.5), Pandas (v2.3.3), dan Matplotlib (v3.7.3). Spesifikasi lengkap lingkungan perangkat lunak disajikan pada Tabel 4.2.
\begin{table}[H]
\centering
\footnotesize
\caption{Spesifikasi Lingkungan Perangkat Lunak}
\label{tab:software-spec}
\begin{tabular}{|c|c|}
\hline
\textbf{Perangkat Lunak / Library} & \textbf{Versi} \\
\hline
Sistem Operasi & Windows 11 (x86\_64) \\
\hline
Python & 3.9.25 \\
\hline
TensorFlow & 2.10.0 \\
\hline
Scikit-learn & 1.6.1 \\
\hline
XGBoost & 2.1.4 \\
\hline
Feature-engine & 1.9.4 \\
\hline
SciPy & 1.13.1 \\
\hline
NumPy & 1.23.5 \\
\hline
Pandas & 2.3.3 \\
\hline
Matplotlib & 3.7.3 \\
\hline
\end{tabular}
\end{table}

\section{Standarisasi Pengujian Eksperimen}
Bagian ini menjelaskan metrik evaluasi yang digunakan untuk menilai performa model klasifikasi serta kualitas data sintetis yang dihasilkan dalam eksperimen. Variabel kontrol diterapkan untuk menjaga konsistensi proses preprocessing, konfigurasi model, dan parameter pelatihan pada seluruh skenario pengujian. Variabel bebas pada penelitian ini adalah penggunaan data sintetis hasil augmentasi berbasis \textit{Variational Autoencoder} (VAE), sedangkan variabel terikat berupa performa klasifikasi dan kualitas data sintetis yang dihasilkan.
\subsection{Konfigurasi Eksperimen}

Bagian ini menjelaskan konfigurasi eksperimen yang digunakan pada proses augmentasi data dan klasifikasi sinyal EEG. Konfigurasi eksperimen meliputi pengaturan dataset, parameter pelatihan model, serta skenario pengujian yang digunakan selama penelitian untuk menjaga konsistensi dan validitas hasil eksperimen.
\subsubsection{Konfigurasi Eksperimen Augmentasi}
\begin{enumerate}
\item \textbf{Dataset:}
    Eksperimen augmentasi menggunakan dataset EEG ASD dari penelitian \cite{eeg_dataset_asd}. Data EEG asli digunakan sebagai data masukan untuk melatih model \textit{Variational Autoencoder} (VAE) dalam menghasilkan sinyal EEG sintetis. Sebelum digunakan pada proses pelatihan model VAE, sinyal EEG melalui tahapan \textit{preprocessing} yang mengacu pada penelitian referensi \cite{paper_ref}, meliputi \textit{filtering}, segmentasi sinyal, serta proses pengurangan artefak untuk meningkatkan kualitas sinyal EEG.

    Eksperimen augmentasi dilakukan menggunakan dua pendekatan berbeda. Pendekatan pertama menggunakan satu model VAE untuk mempelajari sinyal EEG multikanal secara langsung sebelum proses pemisahan kanal. Pendekatan kedua menggunakan model VAE secara terpisah untuk setiap kanal EEG dan setiap kelas data (NORMAL dan ASD), sehingga dihasilkan total 32 model VAE independen.
    \item \textbf{Hyperparameter:}
    Pelatihan model \textit{Variational Autoencoder} (VAE) dilakukan menggunakan ukuran \textit{batch} sebesar 32 selama 70 \textit{epoch}. Dimensi ruang laten (\textit{latent dimension}) yang digunakan sebesar 16 dengan optimizer \textit{Adam} menggunakan \textit{learning rate} sebesar $0.001$, $\beta_1 = 0.5$, dan $\beta_2 = 0.999$. Fungsi \textit{loss} menggunakan kombinasi \textit{Mean Squared Error} (MSE) dan \textit{Kullback-Leibler Divergence} dengan parameter $\beta = 0.5$. Pembagian data dilakukan menggunakan rasio 70\% data pelatihan, 20\% data validasi, dan 10\% data pengujian.

    \item \textbf{Jumlah Pengulangan Eksperimen:}
    Untuk menjamin reliabilitas hasil dan meminimalisir bias akibat inisialisasi acak, setiap eksperimen augmentasi dilakukan secara berulang dengan bobot awal model yang berbeda. Seluruh eksperimen dilakukan sebanyak lima kali pengulangan, kemudian hasil evaluasi akhir diperoleh melalui rata-rata dari seluruh pengulangan eksperimen. Pengulangan eksperimen dilakukan untuk memastikan stabilitas performa model \textit{VAE} dalam menghasilkan data EEG sintetis.
\end{enumerate}

\subsubsection{Konfigurasi Eksperimen Klasifikasi}

\begin{enumerate}
    \item \textbf{Dataset:}
    Eksperimen klasifikasi menggunakan dua skenario dataset, yaitu dataset yang hanya terdiri dari data EEG asli (\textit{real data}) dan dataset hasil augmentasi yang merupakan kombinasi data EEG asli dengan data sintetis hasil generasi model \textit{VAE}. Data sintetis yang dihasilkan pada tahap augmentasi digabungkan kembali menjadi sinyal EEG multikanal, kemudian diproses menggunakan tahapan \textit{preprocessing} yang sama dengan data EEG asli sebagaimana dijelaskan pada penelitian referensi \cite{paper_ref} sebelum digunakan pada proses klasifikasi.

    \item \textbf{Hyperparameter:}
    Eksperimen klasifikasi menggunakan dua model, yaitu model \textit{Deep Neural Network} (DNN) sebagai model pembanding berdasarkan penelitian rujukan serta model \textit{Extreme Gradient Boosting} (XGBoost). Model DNN dilatih menggunakan optimizer Adam dengan \textit{learning rate} sebesar $10^{-3}$, \textit{batch size} 32, dan maksimum 100 epoch. Selain itu, digunakan mekanisme \textit{early stopping} dan \textit{model checkpoint} untuk mencegah \textit{overfitting} selama proses pelatihan. Sementara itu, model XGBoost menggunakan konfigurasi default dari library XGBoost dengan fungsi evaluasi \textit{logloss} serta parameter bawaan library, termasuk jumlah \textit{estimators}. Proses evaluasi model dilakukan menggunakan metode \textit{10-fold Stratified Cross Validation}.

   \item \textbf{Jumlah Pengulangan Eksperimen:} 
    Untuk menjamin reliabilitas hasil dan meminimalisir bias akibat pembagian data maupun inisialisasi model secara acak, setiap eksperimen klasifikasi dilakukan sebanyak lima kali pengulangan. Selain itu, proses evaluasi model dilakukan menggunakan metode \textit{10-fold Stratified Cross Validation}. Seluruh nilai evaluasi yang diperoleh disajikan dalam bentuk rata-rata (\textit{mean}) dan standar deviasi (\textit{standard deviation}) dari sepuluh fold pengujian. 
\end{enumerate}

\subsection{Metrik Evaluasi}

Dalam eksperimen, digunakan sejumlah metrik evaluasi untuk menilai performa model pada proses augmentasi maupun klasifikasi data EEG. Metrik evaluasi digunakan untuk menganalisis kualitas data sintetis yang dihasilkan oleh model \textit{Variational Autoencoder} (VAE) serta mengukur performa model klasifikasi dalam membedakan kelas ASD dan NORMAL.
\subsubsection{Metrik Evaluasi Eksperimen Augmentasi}

Evaluasi pada tahap augmentasi dilakukan untuk menilai kemampuan model \textit{Variational Autoencoder} (VAE) dalam menghasilkan sinyal EEG sintetis yang memiliki karakteristik serupa dengan data EEG asli. Selain memantau perubahan nilai \textit{validation loss} selama proses pelatihan model, evaluasi juga dilakukan menggunakan beberapa metrik kuantitatif yang umum digunakan pada evaluasi model generatif. Metrik evaluasi yang digunakan mengacu pada penelitian \cite{eeg_gan}, meliputi \textit{Frechet Inception Distance} (FID), \textit{Euclidean Distance} (ED), dan \textit{Sliced Wasserstein Distance} (SWD). Ketiga metrik tersebut digunakan untuk mengukur tingkat kemiripan distribusi, karakteristik fitur, serta kedekatan sinyal EEG sintetis terhadap data EEG asli.

\subsubsection{Metrik Evaluasi Eksperimen Klasifikasi}

Evaluasi pada tahap klasifikasi dilakukan untuk menilai performa model dalam membedakan sinyal EEG subjek ASD dan kontrol normal. Evaluasi dilakukan pada dua model klasifikasi, yaitu model \textit{Deep Neural Network} (DNN) sebagai model pembanding berdasarkan penelitian rujukan dan model \textit{Extreme Gradient Boosting} (XGBoost). Kedua model dievaluasi menggunakan skema \textit{10-fold Stratified Cross Validation} untuk menjaga distribusi kelas pada setiap fold pengujian.

Metrik evaluasi yang digunakan meliputi \textit{Precision}, \textit{Recall} (\textit{Sensitivity}), \textit{F1-Score}, \textit{Accuracy}, \textit{Specificity}, serta \textit{Area Under Curve} (AUC). Selain itu, evaluasi juga dilakukan menggunakan \textit{confusion matrix} untuk memperoleh nilai \textit{True Positive} (TP), \textit{True Negative} (TN), \textit{False Positive} (FP), dan \textit{False Negative} (FN) pada setiap fold pengujian. Seluruh metrik evaluasi diterapkan secara konsisten baik pada model DNN maupun model XGBoost untuk menjaga validitas komparasi performa antar model klasifikasi. Seluruh hasil evaluasi disajikan dalam bentuk rata-rata (\textit{mean}) dan standar deviasi (\textit{standard deviation}) dari seluruh fold untuk memberikan gambaran performa model yang lebih stabil dan representatif.

\section{Skenario Pengujian}

Bagian ini menjelaskan skenario pengujian eksperimen yang digunakan untuk mengevaluasi performa klasifikasi EEG dengan augmentasi data berbasis \textit{Variational Autoencoder} (VAE). Eksperimen melibatkan dua skenario utama, yaitu klasifikasi menggunakan data EEG asli (\textit{real data}) dan klasifikasi menggunakan kombinasi data EEG asli dengan data sintetis hasil augmentasi. Selain itu, penelitian juga membandingkan performa model \textit{Deep Neural Network} (DNN) sebagai model \textit{baseline} dan model \textit{Extreme Gradient Boosting} (XGBoost).

\subsection{Eksperimen 1: Evaluasi Kemampuan Model VAE dalam Generasi Sinyal EEG Sintetis}

Eksperimen pertama bertujuan untuk mengevaluasi kemampuan model \textit{Variational Autoencoder} (VAE) dalam mempelajari distribusi data EEG dan menghasilkan sinyal EEG sintetis yang memiliki karakteristik serupa dengan data asli. Eksperimen ini berperan sebagai validasi awal bahwa model VAE mampu merepresentasikan pola sinyal EEG dari masing-masing kelas dan kanal EEG yang digunakan dalam penelitian.

Selain itu, eksperimen ini juga bertujuan untuk menganalisis kualitas data sintetis yang dihasilkan menggunakan beberapa metrik evaluasi data generatif. Hasil dari eksperimen ini digunakan sebagai dasar untuk menentukan apakah data sintetis yang dihasilkan layak digunakan pada tahap augmentasi data untuk proses klasifikasi EEG.

Berdasarkan tujuan tersebut, hipotesis yang diajukan pada eksperimen ini meliputi:

\begin{enumerate}
    \item Model \textit{Variational Autoencoder} (VAE) mampu mencapai konvergensi \textit{loss} selama proses pelatihan pada seluruh kanal EEG dan kelas data yang digunakan.
    
    \item Model VAE mampu menghasilkan sinyal EEG sintetis yang memiliki karakteristik distribusi yang mendekati data EEG asli berdasarkan metrik evaluasi data generatif yang digunakan.
\end{enumerate}

\subsection{Eksperimen 2: Evaluasi Model Klasifikasi Dengan Hanya Menggunakan Sinya EEG Asli}

Eksperimen kedua bertujuan untuk mengevaluasi performa model klasifikasi menggunakan data EEG asli tanpa melibatkan data sintetis hasil augmentasi. Eksperimen ini dilakukan untuk memperoleh nilai performa dasar (\textit{baseline}) yang akan digunakan sebagai acuan pembanding pada eksperimen klasifikasi dengan augmentasi data sintetis.

Eksperimen melibatkan dua model klasifikasi, yaitu model \textit{Deep Neural Network} (DNN) sebagai model \textit{baseline} berdasarkan penelitian rujukan dan model \textit{Extreme Gradient Boosting} (XGBoost). Evaluasi dilakukan menggunakan dataset EEG asli yang telah melalui tahap \textit{preprocessing} sebagaimana dijelaskan pada penelitian referensi.

Berdasarkan tujuan tersebut, hipotesis yang diajukan pada eksperimen ini meliputi:

\begin{enumerate}
    \item Model DNN dan XGBoost mampu melakukan klasifikasi sinyal EEG ASD dan NORMAL menggunakan data EEG asli dengan performa yang stabil.
    
    \item Hasil evaluasi pada eksperimen ini dapat digunakan sebagai nilai \textit{baseline} untuk membandingkan pengaruh penggunaan data sintetis pada eksperimen klasifikasi berikutnya.
\end{enumerate}

\subsection{Eksperimen 3: Evaluasi Model Klasifikasi Dengan Menggunakan Kombinasi Sinyal EEG Asli dan Sinyal EEG Sintetis}

Eksperimen ketiga bertujuan untuk mengevaluasi pengaruh penggunaan data sintetis hasil augmentasi \textit{Variational Autoencoder} (VAE) terhadap performa model klasifikasi EEG. Pada eksperimen ini, proses klasifikasi dilakukan menggunakan kombinasi data EEG asli dan data sintetis yang dihasilkan pada tahap augmentasi.

Eksperimen melibatkan dua model klasifikasi, yaitu model \textit{Deep Neural Network} (DNN) dan model \textit{Extreme Gradient Boosting} (XGBoost). Performa kedua model kemudian dibandingkan dengan hasil pada eksperimen sebelumnya untuk menganalisis pengaruh augmentasi data sintetis terhadap kemampuan klasifikasi sinyal EEG ASD dan NORMAL.

Berdasarkan tujuan tersebut, hipotesis yang diajukan pada eksperimen ini meliputi:

\begin{enumerate}
    \item Penggunaan data sintetis hasil augmentasi berpotensi meningkatkan performa model klasifikasi dibandingkan penggunaan data EEG asli saja.
    
    \item Pengaruh augmentasi data sintetis terhadap performa klasifikasi bergantung pada kemampuan masing-masing model dalam mempelajari pola tambahan dari data hasil augmentasi. Model tertentu dapat memperoleh peningkatan performa yang lebih baik dibandingkan model lainnya seiring bertambahnya jumlah data pelatihan.
\end{enumerate}

\section{Hasil Eksperimen dan Analisis}

Eksperimen dilakukan sesuai dengan skenario pengujian yang telah dijelaskan pada bagian sebelumnya. Bagian ini menyajikan hasil eksperimen serta analisis terhadap performa model pada setiap skenario pengujian, meliputi proses augmentasi data menggunakan \textit{Variational Autoencoder} (VAE) dan proses klasifikasi EEG menggunakan model \textit{Deep Neural Network} (DNN) serta \textit{Extreme Gradient Boosting} (XGBoost). Untuk menjaga fokus pembahasan pada temuan utama penelitian, hanya hasil yang dianggap representatif yang disajikan pada bagian ini, sedangkan keseluruhan hasil eksperimen disertakan pada lampiran.

Hasil eksperimen menunjukkan bahwa penggunaan augmentasi data sintetis memberikan pengaruh yang berbeda terhadap masing-masing model klasifikasi. Berdasarkan hasil evaluasi kualitas data sintetis, model generatif mampu menghasilkan data EEG yang memiliki tingkat kemiripan tertentu terhadap distribusi data asli, meskipun kualitas hasil generasi berbeda pada masing-masing kelas data. Perbedaan tersebut menunjukkan bahwa kualitas data sintetis yang dihasilkan turut dipengaruhi oleh karakteristik dan kualitas distribusi data asli yang digunakan selama proses pelatihan model generatif sehingga kualitas distribusi data EEG tidak sepenuhnya seragam pada seluruh kelas data.

Pada tahap klasifikasi, model \textit{XGBoost} menunjukkan peningkatan performa yang konsisten setelah penambahan data hasil augmentasi. Hal ini mengindikasikan bahwa data sintetis yang dihasilkan masih mampu mempertahankan informasi diskriminatif yang relevan untuk membantu proses klasifikasi ASD. Sebaliknya, model \textit{Deep Neural Network} (DNN) menunjukkan perilaku yang berbeda setelah augmentasi dilakukan. Meskipun beberapa metrik seperti spesifisitas dan presisi mengalami peningkatan, sensitivitas model mengalami penurunan yang cukup signifikan. Hasil ini menunjukkan bahwa model DNN cenderung lebih sensitif terhadap perubahan distribusi data maupun kualitas data sintetis yang dihasilkan oleh model augmentasi.

Perbedaan performa antara kedua model klasifikasi tersebut menunjukkan bahwa efektivitas augmentasi data tidak hanya bergantung pada kualitas data sintetis yang dihasilkan, tetapi juga dipengaruhi oleh karakteristik dan mekanisme pembelajaran dari masing-masing model klasifikasi. Model berbasis \textit{tree boosting} seperti XGBoost cenderung menunjukkan performa yang lebih stabil terhadap variasi distribusi data, sedangkan model berbasis \textit{neural network} lebih sensitif terhadap perubahan distribusi maupun ketidakseimbangan karakteristik data hasil augmentasi.

Walaupun augmentasi data pada penelitian ini belum memberikan peningkatan performa yang konsisten pada seluruh model klasifikasi, hasil eksperimen menunjukkan bahwa pendekatan generatif dalam augmentasi memiliki potensi untuk meningkatkan kemampuan generalisasi model klasifikasi EEG pada kasus ASD. Pembahasan lebih mendalam mengenai hasil dari masing-masing eksperimen dijelaskan pada sub-bagian berikut.

\subsection{Eksperimen 1: Evaluasi Kemampuan Model VAE dalam Generasi Sinyal EEG Sintetis}
\begin{table}[H]
\centering
\footnotesize
\caption{Hasil Evaluasi Kualitas Data Sintetis EEG pada Kelas Normal}
\label{tab:normal_generation_results}
\begin{tabular}{|c|c|c|c|}
\hline
Channel & ED Norm & FID & SWD \\
\hline
Fp1 & -0.148 & 91.313 & 0.264 \\
\hline
Fp2 & 0.019 & 60.794 & 0.213 \\
\hline
F3  & 0.086 & 109.476 & 0.199 \\
\hline
F4  & 0.037 & 86.644 & 0.172 \\
\hline
C3  & 0.127 & 110.533 & 0.184 \\
\hline
C4  & 0.058 & 95.488 & 0.145 \\
\hline
P3  & 0.049 & 129.378 & 0.132 \\
\hline
P4  & 0.145 & 132.878 & 0.147 \\
\hline
O1  & 0.183 & 115.336 & 0.202 \\
\hline
O2  & 0.222 & 118.257 & 0.186 \\
\hline
F7  & 0.260 & 121.811 & 0.181 \\
\hline
F8  & 0.271 & 114.876 & 0.169 \\
\hline
T3  & 0.193 & 152.974 & 0.136 \\
\hline
T4  & 0.393 & 180.492 & 0.233 \\
\hline
T5  & 0.588 & 170.312 & 0.213 \\
\hline
T6  & 0.531 & 164.081 & 0.235 \\
\hline
\end{tabular}
\end{table}
\begin{table}[H]
\centering
\footnotesize
\caption{Hasil Evaluasi Kualitas Data Sintetis EEG pada Kelas ASD}
\label{tab:asd_generation_results}
\begin{tabular}{|c|c|c|c|}
\hline
Channel & ED Norm & FID & SWD \\
\hline
Fp1 & 0.003 & 34.857 & 0.136 \\
\hline
Fp2 & 0.012 & 26.858 & 0.111 \\
\hline
F3  & 0.047 & 49.890 & 0.153 \\
\hline
F4  & 0.021 & 45.219 & 0.104 \\
\hline
C3  & -0.020 & 42.832 & 0.109 \\
\hline
C4  & 0.014 & 49.008 & 0.118 \\
\hline
P3  & -0.026 & 49.183 & 0.113 \\
\hline
P4  & -0.029 & 48.179 & 0.084 \\
\hline
O1  & 0.014 & 38.777 & 0.118 \\
\hline
O2  & 0.038 & 33.451 & 0.105 \\
\hline
F7  & 0.056 & 32.921 & 0.101 \\
\hline
F8  & 0.052 & 31.354 & 0.107 \\
\hline
T3  & 0.081 & 48.797 & 0.110 \\
\hline
T4  & 0.070 & 50.006 & 0.090 \\
\hline
T5  & 0.149 & 52.433 & 0.110 \\
\hline
T6  & 0.120 & 51.330 & 0.108 \\
\hline
\end{tabular}
\end{table}
Eksperimen pertama bertujuan untuk mengevaluasi kemampuan model \textit{Variational Autoencoder} (VAE) dalam mempelajari distribusi sinyal EEG dan menghasilkan data sintetis yang menyerupai data asli pada masing-masing kanal EEG sebelum digunakan pada tahap klasifikasi ASD.

Berdasarkan hasil evaluasi pada Tabel \ref{tab:normal_generation_results} dan Tabel \ref{tab:asd_generation_results}, model VAE mampu menghasilkan sinyal sintetis yang mempertahankan pola umum sinyal EEG asli pada kedua kelas data. Kemiripan tersebut juga terlihat pada visualisasi domain waktu dan domain frekuensi pada Gambar X.X dan Gambar X.X, di mana sinyal sintetis mampu mengikuti pola amplitudo utama serta distribusi spektral frekuensi dominan dari sinyal asli.

Secara umum, kualitas data sintetis pada kelas ASD menunjukkan hasil yang lebih baik dan lebih stabil dibandingkan kelas Normal. Hal ini ditunjukkan oleh nilai \textit{Fréchet Inception Distance} (FID), \textit{Sliced Wasserstein Distance} (SWD), dan normalisasi \textit{Euclidean Distance} (ED) yang relatif lebih rendah pada sebagian besar kanal EEG kelas ASD. Kanal frontal seperti Fp2, F7, dan F8 pada kelas ASD menghasilkan nilai FID yang rendah, menunjukkan tingkat kemiripan distribusi yang tinggi terhadap data asli.

Sebaliknya, beberapa kanal pada kelas Normal, khususnya kanal temporal seperti T4, T5, dan T6, menunjukkan nilai FID dan normalisasi ED yang lebih tinggi dibandingkan kanal lainnya. Secara empiris, hal ini mengindikasikan bahwa distribusi sinyal pada kanal tersebut lebih sulit dipelajari oleh model generatif, kemungkinan akibat variasi pola sinyal dan kompleksitas distribusi yang lebih tinggi.

Perbedaan kualitas hasil generasi juga terlihat selama proses pelatihan model. Berdasarkan grafik evolusi \textit{validation loss} pada Gambar X.X, model pada data ASD cenderung mencapai konvergensi yang lebih stabil dengan nilai loss yang lebih rendah, sedangkan model pada data Normal menunjukkan \textit{validation loss} yang lebih tinggi dan mengalami plateau lebih awal.

Meskipun kualitas hasil generasi tidak sepenuhnya seragam pada seluruh kanal, hasil eksperimen menunjukkan bahwa model VAE tetap mampu mempertahankan karakteristik utama sinyal EEG pada domain waktu maupun domain frekuensi. Temuan ini menunjukkan bahwa pendekatan augmentasi berbasis model generatif memiliki potensi untuk meningkatkan jumlah dan variasi data EEG pada proses klasifikasi ASD.

\subsection{Eksperimen 2: Evaluasi Model Klasifikasi Menggunakan Data EEG Asli}

Eksperimen kedua bertujuan untuk mengevaluasi performa model klasifikasi menggunakan data EEG asli tanpa melibatkan data sintetis hasil augmentasi. Eksperimen ini dilakukan untuk memperoleh nilai performa dasar (\textit{baseline}) yang akan digunakan sebagai acuan pembanding pada eksperimen klasifikasi dengan augmentasi data sintetis.

Eksperimen melibatkan dua model klasifikasi, yaitu model \textit{Deep Neural Network} (DNN) sebagai model \textit{baseline} berdasarkan penelitian rujukan dan model \textit{Extreme Gradient Boosting} (XGBoost). Seluruh pengujian dilakukan menggunakan data EEG asli yang telah melalui tahap \textit{preprocessing} sebagaimana dijelaskan pada penelitian referensi. Evaluasi performa dilakukan menggunakan metrik akurasi, sensitivitas, spesifisitas, presisi, \textit{F1-score}, dan \textit{Area Under Curve} (AUC). Hasil eksperimen ditunjukkan pada Tabel X.X.

Berdasarkan hasil evaluasi, kedua model mampu melakukan klasifikasi sinyal EEG ASD dan Normal menggunakan data EEG asli. Namun, terdapat perbedaan performa yang cukup signifikan antara model DNN dan XGBoost. Model XGBoost secara konsisten menunjukkan performa yang lebih tinggi pada hampir seluruh metrik evaluasi dibandingkan model DNN. Nilai sensitivitas, spesifisitas, presisi, \textit{F1-score}, dan AUC yang diperoleh XGBoost menunjukkan bahwa model mampu melakukan pemisahan kelas dengan sangat baik dan stabil pada sebagian besar fold pengujian.

Sebaliknya, model DNN menunjukkan performa yang relatif lebih rendah dengan variasi hasil yang lebih besar antar fold pengujian. Hal ini terlihat dari nilai standar deviasi yang cukup tinggi pada beberapa metrik evaluasi, yang mengindikasikan bahwa performa model belum sepenuhnya stabil. Selain itu, jumlah \textit{false positive} dan \textit{false negative} pada model DNN juga relatif lebih tinggi dibandingkan XGBoost, sehingga kemampuan generalisasi model terhadap data EEG masih terbatas.

Perbedaan performa tersebut menunjukkan bahwa karakteristik data EEG pada penelitian ini cenderung lebih efektif dipelajari oleh model berbasis \textit{gradient boosting} dibandingkan model DNN. Secara empiris, XGBoost memiliki kemampuan yang lebih baik dalam menangani jumlah data yang relatif terbatas serta hubungan fitur yang kompleks, sehingga menghasilkan performa klasifikasi yang lebih stabil dibandingkan DNN pada konfigurasi pengujian yang digunakan.

Secara keseluruhan, hasil eksperimen ini menunjukkan bahwa model DNN dan XGBoost memiliki karakteristik performa yang berbeda dalam melakukan klasifikasi sinyal EEG ASD menggunakan data asli. Oleh karena itu, kedua model selanjutnya digunakan pada eksperimen berikutnya untuk mengevaluasi bagaimana penggunaan data sintetis hasil augmentasi memengaruhi performa dan kemampuan generalisasi masing-masing model klasifikasi.

\subsection{Eksperimen 3: Evaluasi Pengaruh Augmentasi Data EEG terhadap Performa Model Klasifikasi}

Hasil eksperimen menunjukkan bahwa penggunaan data sintetis memberikan pengaruh yang berbeda pada masing-masing model klasifikasi. Secara umum, augmentasi data memberikan peningkatan performa yang signifikan pada model XGBoost, namun menghasilkan perubahan karakteristik performa yang berbeda pada model DNN.

Berdasarkan hasil evaluasi pada Tabel X.X, model XGBoost mengalami peningkatan performa pada hampir seluruh metrik evaluasi setelah penambahan data sintetis. Nilai sensitivitas meningkat dari 0.971 menjadi 0.987, sedangkan spesifisitas meningkat hingga mendekati sempurna yaitu dari 0.978 menjadi 0.999. Selain itu, akurasi model juga meningkat dari 0.974 menjadi 0.995 dengan nilai presisi yang sangat tinggi. Hasil tersebut menunjukkan bahwa data sintetis yang dihasilkan mampu membantu model XGBoost mempelajari distribusi data EEG dengan lebih baik serta meningkatkan kemampuan generalisasi model terhadap kedua kelas data.

Sebaliknya, model DNN menunjukkan perubahan performa yang tidak sepenuhnya konsisten setelah penggunaan data sintetis. Meskipun akurasi meningkat dari 0.585 menjadi 0.724 dan nilai presisi meningkat secara signifikan dari 0.537 menjadi 0.833, nilai sensitivitas justru mengalami penurunan yang cukup besar dari 0.630 menjadi 0.212. Sementara itu, spesifisitas meningkat secara signifikan dari 0.539 menjadi 0.967. Pola ini menunjukkan bahwa model DNN cenderung menjadi lebih konservatif dalam melakukan prediksi kelas ASD setelah penambahan data sintetis.

Secara empiris, kondisi tersebut mengindikasikan bahwa model DNN lebih sering memprediksi kelas Normal dibandingkan kelas ASD setelah proses augmentasi dilakukan. Akibatnya, jumlah \textit{false positive} berhasil ditekan sehingga nilai presisi dan spesifisitas meningkat, namun di sisi lain jumlah \textit{false negative} meningkat sehingga sensitivitas mengalami penurunan. Hal ini menunjukkan bahwa penambahan data sintetis tidak selalu meningkatkan kemampuan model DNN dalam mengenali pola ASD secara seimbang.

Perbedaan respons kedua model terhadap augmentasi kemungkinan berkaitan dengan karakteristik arsitektur masing-masing model. Model XGBoost cenderung lebih robust terhadap perubahan distribusi data dan lebih efektif dalam memanfaatkan tambahan variasi fitur dari data sintetis. Sebaliknya, model DNN lebih sensitif terhadap kualitas dan distribusi data augmentasi, terutama ketika kualitas data sintetis antar kelas tidak sepenuhnya seimbang. Pada eksperimen sebelumnya, kualitas data sintetis pada kelas Normal menunjukkan deviasi distribusi yang lebih besar dibandingkan kelas ASD. Kondisi ini kemungkinan turut memengaruhi proses pembelajaran model DNN sehingga distribusi keputusan model menjadi kurang seimbang setelah augmentasi dilakukan.

Meskipun demikian, hasil eksperimen menunjukkan bahwa augmentasi berbasis model generatif tetap memiliki potensi dalam meningkatkan performa klasifikasi EEG, terutama pada model berbasis \textit{gradient boosting} seperti XGBoost. Temuan ini juga menunjukkan bahwa efektivitas data sintetis sangat dipengaruhi oleh karakteristik model klasifikasi yang digunakan serta kualitas distribusi data sintetis yang dihasilkan oleh model augmentasi.